\section{Min Stack}
\label{sec:sq-min-stack}

\problemstatement{Design a stack supporting $\textit{push}(x)$,
$\textit{pop}()$, $\textit{top}()$, and $\textit{getMin}()$ -- the last
returning the minimum element currently in the stack -- all in $O(1)$
time.}

\begin{patternbox}
\emph{``Maintain a running aggregate (min, max, sum) that must stay
correct as elements are pushed \emph{and} popped''} is the keyword
pattern here. A plain running variable breaks the moment a \textit{pop}
removes the element that was responsible for the current minimum -- there
would be no way to recover the \emph{previous} minimum without
re-scanning. The fix is to keep a full \textbf{history} of the minimum, one
entry per stack position, so that popping the main stack and popping the
history together always keeps both in sync.
\end{patternbox}

\subsection*{Understanding the problem carefully}

The difficulty is specifically about \textit{pop}: tracking the minimum
as elements are pushed is trivial (compare the new element to the current
minimum). The problem is that after a \textit{pop}, the minimum might need
to \emph{revert} to whatever it was before the popped element was pushed
-- and that historical value must be available instantly, not
recomputed.

\begin{ideabox}[A Second Stack Tracks the Minimum-So-Far at Every Depth]
Maintain an auxiliary \textit{minStack}, parallel to the main stack. On
$\textit{push}(x)$: push $x$ onto the main stack as usual, and push
$\min(x, \textit{minStack.top}())$ onto \textit{minStack} (or just $x$ if
\textit{minStack} is currently empty). On $\textit{pop}()$: pop both
stacks together. $\textit{getMin}()$ simply reads \textit{minStack}'s
top.
\end{ideabox}

\subsection*{Why popping both stacks together keeps the invariant correct}

\textit{minStack}'s $i$-th entry (counting from the bottom) is, by
construction, the minimum of the main stack's bottom $i$ elements at the
moment the $i$-th element was pushed -- and since both stacks always have
the same height and are pushed to and popped from in lockstep,
\textit{minStack}'s current top is always exactly the minimum of
\emph{whatever the main stack currently contains}, not some stale value
from before a pop. This is the same ``parallel history'' idea that makes
undo/redo systems work: rather than computing an aggregate forward only,
store enough history to reconstruct any previous state instantly.

\subsection*{Algorithm}

\begin{enumerate}
  \item $\textit{push}(x)$: push $x$ onto \textit{mainStack}. Push
        $\min(x, \text{current } \textit{getMin}())$ onto
        \textit{minStack} (using $x$ itself if \textit{minStack} was
        empty).
  \item $\textit{pop}()$: pop both \textit{mainStack} and
        \textit{minStack}.
  \item $\textit{top}()$: return \textit{mainStack}'s top.
  \item $\textit{getMin}()$: return \textit{minStack}'s top.
\end{enumerate}

\subsection*{Dry run}

\dryrun{Push $5, 3, 7, 2$; pop once; query getMin.}

\begin{itemize}
  \item Push $5$: main$=[5]$, min$=[5]$.
  \item Push $3$: main$=[5,3]$, min$=[5,\min(3,5){=}3]$.
  \item Push $7$: main$=[5,3,7]$, min$=[5,3,\min(7,3){=}3]$.
  \item Push $2$: main$=[5,3,7,2]$, min$=[5,3,3,\min(2,3){=}2]$.
  \item Pop: both stacks lose their top. main$=[5,3,7]$, min$=[5,3,3]$.
  \item getMin(): top of min stack is $3$.
\end{itemize}

Correct: after removing $2$, the remaining stack $[5,3,7]$ has minimum
$3$.

\subsection*{C++ implementation}

\begin{lstlisting}
#include <bits/stdc++.h>
using namespace std;

class MinStack {
    stack<int> mainStack, minStack;

public:
    void push(int x) {
        mainStack.push(x);
        if (minStack.empty()) minStack.push(x);
        else minStack.push(min(x, minStack.top()));
    }

    void pop() {
        mainStack.pop();
        minStack.pop();
    }

    int top() {
        return mainStack.top();
    }

    int getMin() {
        return minStack.top();
    }
};

int main() {
    MinStack s;
    s.push(5); s.push(3); s.push(7); s.push(2);
    s.pop();
    cout << "Min: " << s.getMin() << endl;
    return 0;
}
\end{lstlisting}

\begin{complexitybox}
\textbf{Time Complexity.} $O(1)$ for every operation.
\textbf{Space Complexity.} $O(n)$, doubled from a plain stack since every
push adds one entry to \emph{both} stacks. This can be reduced to
$O(1)$ \emph{extra} space (still $O(n)$ total, but without a second full
stack) using a trick where only the previous minimum is pushed onto the
main stack itself when a new minimum is encountered, encoded so it can be
distinguished and restored on pop -- a worthwhile optimization to know
exists, though the two-stack version above is the version to derive and
explain first.
\end{complexitybox}

\begin{takeawaybox}
Whenever an aggregate (min, max, or anything else that a plain running
variable cannot correctly \emph{revert} after a removal) must stay valid
across both insertions and removals, keep a parallel history stack rather
than trying to recompute the aggregate from scratch on every pop.
\end{takeawaybox}
